\documentclass[12pt,twoside]{article}

\usepackage[english]{babel}
\usepackage[utf8]{inputenc}
\usepackage{sansmathfonts}
\usepackage[T1]{fontenc}
\renewcommand*\familydefault{\sfdefault} %% Only if the base font of the document is to be sans serif
%\usepackage{tabularx}
\usepackage{dashrule}
\usepackage{tikz}
\usepackage{tabularray}
\usepackage[export]{adjustbox}
%\captionsetup[longtable]{skip=30pt}
\usepackage{arydshln}
\usepackage{ctable}
\usepackage{colortbl}
\usepackage{multicol,lipsum}
\usepackage{amsmath}
\usepackage{amssymb,amsthm,amsmath,dsfont,mathrsfs}
\usepackage[colorinlistoftodos]{todonotes}
\usepackage{geometry}
\geometry{
a4paper,
total={170mm,257mm},
left=15mm,
right=15mm,
top=15mm,
bottom=10mm,
}
\usepackage{natbib}
\usepackage{soul}
\usepackage{setspace}
\usepackage{hyperref}
\usepackage{xcolor}
\usepackage{cancel}
\color{darkgray}
\usepackage{graphicx}
\usepackage{fancyhdr}
\pagestyle{fancy}
\renewcommand{\headrulewidth}{0pt}


\newcommand{\T}[1]{\textbf{\large #1} \newline}

\newcommand{\TM}[1]{ \textbf{ \newline #1}  }
\setlength\columnsep{30pt}

\definecolor{bluebell}{rgb}{0.39, 0.44, 1}
\definecolor{laranja}{rgb}{1, 0.58, 0.33}
\definecolor{verde}{rgb}{0, 0.46, 0.37}
\definecolor{babypink}{rgb}{0.85, 0.65, 0.65} 
\definecolor{cinza}{rgb}{0.3, 0.3, 0.3} 
\definecolor{iris}{rgb}{0.39, 0.44, 1} 
\definecolor{babyblue}{rgb}{0.13, 0.59, 0.95} 

\color{cinza} 
 
\newcommand{\C}[1]{\textcolor{cinza}{#1}}
\newcommand{\BI}[1]{\textbf{\textcolor{iris}{#1}}}
\newcommand{\BL}[1]{\textbf{\textcolor{laranja}{#1}}}
\newcommand{\BV}[1]{\textbf{\textcolor{verde}{#1}}}
\newcommand{\BC}[1]{\textbf{\textcolor{cinza}{#1}}}
\newcommand{\BB}[1]{\textbf{\textcolor{babyblue}{#1}}}

\begin{document}
\pagenumbering{gobble}
\setstretch{1}

\setstretch{1} 
\normalsize 
\setlength{\tabcolsep}{0em} 
\setlength\parindent{0pt} 
\begin{tabular}{p{0.19\textwidth} p{0.61\textwidth} p{0.2\textwidth}}
\BB{\Large Exercícios } & \footnotesize \textcolor{iris}{Módulo Básico} \BL{-} \textcolor{verde}{Conjuntos} \BL{-} Nível I & Nota:    \\  
\end{tabular}
\vspace{0.3cm}

\small
\textcolor{laranja}{Instruções: } \textcolor{cinza}{asd}
\vspace{0.1cm}

\setstretch{1.3}
\footnotesize
\begin{tblr}{
colspec={p{0.46\textwidth}|[2pt,laranja]p{0.46\textwidth}},
columns = {colsep=0pt, rightsep=10pt},
column{2} = {leftsep=8pt},
rows={valign=h, ht=8.200000000000001cm},
}
\BI{1)~}Qual a soma dos elementos da intersecção \(A \cap B ~\)?\vspace{5px}  \newline$A=\{5,13,15,18,19,21,24\}\newline B=\{1,6,9,13,18,19,24\}$ &\BI{2)~}Quantos elementos tem a união \(A \cup B ~\)?\vspace{5px}  \newline$A=\{4,10,19,22,23\}\newline B=\{9,10,18,23\}$ \\ 
\BI{3)~}Quantos elementos tem a união \(A \cup B ~\)?\vspace{5px}  \newline$A=\{0,1,3,6,7,8,10,12,14,19\}\newline B=\{3,5,6,8,16,18,24\}$ &\BI{4)~}Qual o menor elemento da união \(A \cup B ~\)?\vspace{5px}  \newline$A=\{5,7,8,14,22\}\newline B=\{0,5,14,18,19,21\}$ \\ 
\BI{5)~}Qual o menor elemento da intersecção \(A \cap B ~\)?\vspace{5px}  \newline$A=\{2,4,6,16,20,21\}\newline B=\{4,6,11,18,20\}$ &\BI{6)~}Qual o menor elemento de \(A - B ~\)?\vspace{5px}  \newline$A=\{4,6,7,8,12,24\}\newline B=\{8,12,15,24\}$ \\ 
\end{tblr}

\newpage 
\setstretch{1} 
\normalsize 
\setlength{\tabcolsep}{0em} 
\setlength\parindent{0pt} 
\BB{\Large Resolução } \quad \footnotesize \textcolor{iris}{Módulo Básico} \BL{-} \textcolor{verde}{Conjuntos} \BL{-} Nível I
\vspace{0.3cm}
\small

\textcolor{laranja}{Instruções: } \textcolor{cinza}{asd}
\vspace{0.1cm}



\setstretch{1.5}
\footnotesize
\begin{tblr}{
colspec={p{0.46\textwidth}|[1pt,cinza]p{0.46\textwidth}},
columns = {colsep=0pt, rightsep=10pt},
column{2} = {leftsep=8pt},
rows={valign=h, ht=8.200000000000001cm},
}
\BI{1)~}{Qual a soma dos elementos da intersecção \(A \cap B ~\)?}\vspace{5px}  \newline 
{ $A=\{5,13,15,18,19,21,24\}\newline B=\{1,6,9,13,18,19,24\}$} \vspace{5px} \newline 
{$A \cap B = \{13,18,19,24\} \\\text{Soma dos elementos de}~ A \cap B \text{ é: } \\13+18+19+24=74$} &\BI{2)~}{Quantos elementos tem a união \(A \cup B ~\)?}\vspace{5px}  \newline 
{ $A=\{4,10,19,22,23\}\newline B=\{9,10,18,23\}$} \vspace{5px} \newline 
{$A \cup B = \{4,9,10,18,19,22,23\} \\|A \cup B|=7$} \\ 
\BI{3)~}{Quantos elementos tem a união \(A \cup B ~\)?}\vspace{5px}  \newline 
{ $A=\{0,1,3,6,7,8,10,12,14,19\}\newline B=\{3,5,6,8,16,18,24\}$} \vspace{5px} \newline 
{$A \cup B = \{0,1,3,5,6,7,8,10,12,14,16,18,19,24\} \\|A \cup B|=14$} &\BI{4)~}{Qual o menor elemento da união \(A \cup B ~\)?}\vspace{5px}  \newline 
{ $A=\{5,7,8,14,22\}\newline B=\{0,5,14,18,19,21\}$} \vspace{5px} \newline 
{$A \cup B = \{0,5,7,8,14,18,19,21,22\} \\\text{Menor elemento de}~ A \cup B \text{ é }0$} \\ 
\BI{5)~}{Qual o menor elemento da intersecção \(A \cap B ~\)?}\vspace{5px}  \newline 
{ $A=\{2,4,6,16,20,21\}\newline B=\{4,6,11,18,20\}$} \vspace{5px} \newline 
{$A \cap B = \{4,6,20\} \\\text{Menor elemento de}~ A \cap B \text{ é }4$} &\BI{6)~}{Qual o menor elemento de \(A - B ~\)?}\vspace{5px}  \newline 
{ $A=\{4,6,7,8,12,24\}\newline B=\{8,12,15,24\}$} \vspace{5px} \newline 
{$A - B = \{4,6,7\} \\\text{Menor elemento de}~ A - B \text{ é }4$} \\ 
\end{tblr}

\vfill
\begin{flushright}
\setlength{\tabcolsep}{0em}
\setlength\parindent{0pt}
\footnotesize \textbf{pratiqueja.com}
\end{flushright}

\end{document}
